\begin{table}[ht]
\centering
\caption{Admissibility-violation taxonomy for unconstrained (C0) candidates on
model-confirmed Class~3 patients, resolving the apparent tension between the
candidate-level and patient-level statistics. \emph{Candidate-level} rates use
all generated C0 candidates as the denominator; \emph{patient-level} uses
patients. Categories: \emph{direction} (an anthropometric increases, or a
reduce-only nutrient increases), \emph{floor} (a value falls below an absolute
Layer-1 safety floor), \emph{coupling} (BMI/waist/weight move off the locked
constant-height ratio), \emph{range} (any other permitted-band breach). Nearly
every C0 candidate breaches at least one constraint, which is why a post-hoc
filter requiring a \emph{fully} admissible candidate removes essentially all of
them, even though the milder patient-level ``any violation'' rate is what
appears in Table~\ref{tab:quality_revised}.}
\label{tab:violation_taxonomy}
\begin{tabular}{lc}
\toprule
Quantity & Value \\
\midrule
  Candidate-level: mean fraction of C0 candidates violating & 100.0\% \\
  \quad direction (wrong-sign move) & 92.2\% \\
  \quad floor (below safety floor) & 37.8\% \\
  \quad coupling (anthropometric ratio broken) & 80.0\% \\
  \quad range (other band breach) & 0.8\% \\
  Patient-level: patients with $\ge$1 violating candidate & 100.0\% \\
\bottomrule
\end{tabular}
\end{table}
